\documentclass[10pt,a4paper]{amsart}
\usepackage[margin=19mm]{geometry}
\usepackage{amsmath,amssymb,mathtools}
\usepackage[scr=rsfs,cal=euler]{mathalfa}
\usepackage{xfrac}
\usepackage[unicode,hidelinks,bookmarks=false]{hyperref}
\newcommand{\ie}{i.e.\ }
\newcommand{\cf}{cf.\ }
\newcommand{\resp}{respectively\ }
\newcommand{\Subsection}{\subsection}
\providecommand{\nobreakdash}{\nobreak\hbox{-}}
\setlength{\emergencystretch}{2em}
\multlinegap0pt
\usepackage{bm}
\usepackage{bbm}
\usepackage{dsfont}
\usepackage{xspace}
\usepackage{cancel}
\usepackage{mathtools}
\usepackage[shortlabels]{enumitem}
\usepackage{amsthm}
\makeatletter
\@ifundefined{restatable}{\newtheorem{restatable}{Restatable}}{}
\makeatother
\title[Situationally-Aware Dynamics Learning]{Situationally-Aware Dynamics Learning}
\author{Alejandro Murillo-Gonzalez, Lantao Liu}
\date{Source: arXiv:2505.19574v3 (CC BY 4.0)}
\begin{document}
\maketitle

\noindent\emph{Source and licence.} Situationally-Aware Dynamics Learning. Version arXiv:2505.19574v3 (CC BY 4.0), distributed under the Creative Commons Attribution 4.0 International licence, as recorded in the arXiv licence metadata for this exact version and shown on the abstract page https://arxiv.org/abs/2505.19574v3. Full rights notice: licenses/arxiv-2505.19574-CC-BY-4.0.txt.\par\smallskip

\noindent\textit{Editorial scope.} Counted unit: the Lemma whose source label is \texttt{sitproblemma} in the article (source0.tex lines 412--426, source label \texttt{sitproblemma}), delivered together with the article's own complete original proof of it (source0.tex lines 428--460, one \texttt{proof} environment of 33 lines). No mathematics is written, repaired or re-derived: statement and proof are the article's own text line for line. Required context retained uncounted: eq:growth-prob-def (lines 277--279); eq:changepoint-prob-def (lines 281--283); eq:sampling-model (lines 389--391). Every definition, display and label of those lines is retained unchanged. Omitted from this package and not redistributed: 1--276, 280--280, 284--388, 392--411, 427--427, 461--1170. Nothing inside a retained excerpt is omitted or altered: every retained body is delivered exactly as the article's lines are written, so no omission receipt of any kind accompanies this package. Citations: the retained excerpts carry no citation command at all, so no bibliography entry of the article is transcribed and no reference is redistributed. Reference adaptation: none was needed and none was made; every reference inside the delivered unit resolves inside it, so no unresolved reference or citation is produced. Printed slips kept literal: no printed word, symbol, hypothesis or number of the counted statement or of its proof was altered. Layout and class: the article's own class is \texttt{sagej} with packages including hyperref, multicol, graphicx, epstopdf, mathtools, amsfonts, algorithm, algpseudocode, algorithmicx, wrapfig, subfiles, booktabs, flushend, amsmath; the wrapper is the lane's pinned \texttt{amsart} editorial wrapper at 10pt on A4, which restates the theorem-like environments the retained text needs and adds the article's own macro definitions that the retained text needs (none), with extra wrapper packages (bm, bbm, dsfont, xspace, cancel, mathtools, enumitem). The delivered numbering is the unit's own. Assets: no retained span contains a figure environment, a graphics-inclusion command, a PDF-page inclusion, a table or a TikZ picture; no image file is copied into this package, the package declares no asset dependency and no image file is redistributed with it. Licence: version arXiv:2505.19574v3 of this article is distributed under the Creative Commons Attribution 4.0 International licence, as recorded in the arXiv licence metadata for this exact version and shown on the abstract page https://arxiv.org/abs/2505.19574v3. A full rights notice is delivered at licenses/arxiv-2505.19574-CC-BY-4.0.txt. Characters: every retained excerpt is pure ASCII, so the delivered engine, XeLaTeX with its default Latin Modern text font, reports no missing character.
\begin{align} \label{eq:growth-prob-def}
    P(r_t = r_{t-1}+1, &\boldsymbol{x}_{1:t}) =\\ &P(r_{t-1}, \boldsymbol{x}_{1:t-1}) \pi_t^{(r)} (1 - H(r_{t-1})). \nonumber
\end{align} 

\begin{align} \label{eq:changepoint-prob-def}
    P(r_t=0, \boldsymbol{x}_{1:t}) = \sum_{\substack{r_{t-1}}} P(r_{t-1}, \boldsymbol{x}_{1:t-1}) \pi_t^{(r)} H(r_{t-1}).
\end{align}

\begin{align} \label{eq:sampling-model}
    p(\mathcal{D} | \boldsymbol{\mu}, \Lambda) &= \prod_{\substack{\boldsymbol{x}_i \in \mathcal{D}}} \mathcal{N} (\boldsymbol{x}_i | \boldsymbol{\mu}, \Lambda)  \\&= \left( \frac{|\Lambda|}{(2 \pi)^{d}} \right)^{n/2} \nonumber \\ & \quad\quad \cdot \exp{ \left( -\frac{1}{2} \sum_{\substack{i=1}}^n (\boldsymbol{x}_i - \boldsymbol{\mu})^\top \Lambda (\boldsymbol{x}_i - \boldsymbol{\mu}) \right) \nonumber}. 
\end{align}

\begin{restatable}[Growth and Changepoint Probabilities]{lemma}{sitproblemma} \label{lemma:sit-prob}
        Let $\boldsymbol{x}_t$ be the observation received at time $t$. Let $\boldsymbol{\eta}^{(r)} = \{\boldsymbol{\mu}_{i}^{(r)}, T_{i}^{(r)}, \nu_{i}^{(r)}, \kappa_{i}^{(r)}\}_{i=t_{cp}}^{t}$ be the parameters of the UDGP $(r)$ that started at time $t_{cp} < t$. Then, 
        \begin{itemize}[leftmargin=0.45cm]
            \item The \textbf{Growth Probability}, or probability that we stayed in the same UDGP, at time $t$ is: %\vspace{-1mm}
            \begin{align} \label{eq:permanence-prob}
                P(r_t = &~r_{t-1} + 1, \boldsymbol{x}_{t_{cp}:t}) = \\ &\frac{\lambda - 1}{\lambda} P(r_{t-1}, \boldsymbol{x}_{t_{cp}:t-1})  
                \mathcal{N} (\boldsymbol{x}_t | \boldsymbol{\mu}_{t-1}^{(r)}, \Lambda_{t-1}^{(r)}). \nonumber
            \end{align}
            \item The \textbf{Changepoint Probability}, or probability that we moved to a different UDGP, at time $t$ is:
            \begin{align} \label{eq:jump-prob}
                P(r_t = 0, &~\boldsymbol{x}_{t_{cp}:t}) = \\ & \frac{1}{\lambda} \sum_{\substack{\tau = 0}}^{r_{t-1}} P(\tau, \boldsymbol{x}_{t_{cp}:t-1}) 
                \mathcal{N} (\boldsymbol{x}_t | \boldsymbol{\mu}_{\tau}^{(r)}, \Lambda_{\tau}^{(r)}). \nonumber
            \end{align}
        \end{itemize}
\end{restatable}

\begin{proof} We assume that the situation the robot starts in begins at the exact moment the robot processes its first observation. Therefore, $P(r_0 = 0, \boldsymbol{x}_{0:0}) = 1$. This becomes the initial condition for the recursive computation key in the BOCD framework to find the distribution of the run length $P(r_t | \boldsymbol{x}_{1:t})$. Then, we can obtain the probability of staying in the current UDGP or jumping to a new UDGP as follows:

    \begin{itemize}[leftmargin=0.45cm]
        \item \textit{\textbf{Growth Probability:}} As defined in Eq. \eqref{eq:growth-prob-def}, the probability that the current run length increases is $P(r_t = r_{t-1}+1, \boldsymbol{x}_{1:t}) = P(r_{t-1}, \boldsymbol{x}_{1:t-1}) \pi_t^{(r)} (1 - H(r_{t-1}))$. For us, $\pi_t^{(r)}$ is the probability of the data given by Eq. \eqref{eq:sampling-model}, our sampling model. We also established before that $H(r_{t-1}) = 1 / \lambda$. Then, the growth probability at $t=1$ is:
        \begin{align*}
            P(r_1 = 1,&~\boldsymbol{x}_{0:1} = \{\boldsymbol{x}_1\}) \\&= P(r_0, \boldsymbol{x}_{0:0}) \cdot \pi_1^{(r)} \cdot (1 - H(0)) \nonumber \\
            &= 1 \cdot P(\{\boldsymbol{x}_1\} | \boldsymbol{\eta}_0) \cdot (1 - \frac{1}{\lambda}) \nonumber \\
            &= \frac{\lambda - 1}{\lambda} \mathcal{N}(\boldsymbol{x}_1 | \boldsymbol{\mu}_0, \Lambda_0). \nonumber
        \end{align*}

        And the growth probability at $t > 1$ is:
        \begin{align*}
            P(&r_t=r_{t-1}+1, \boldsymbol{x}_{0:t}) \\&= P(r_{t-1}, \boldsymbol{x}_{0:t-1}) \cdot \pi_t^{(r)} \cdot (1 - H(r_{t-1})) \nonumber \\
            &= P(r_{t-1}, \boldsymbol{x}_{0:t-1}) \cdot \prod_{\substack{r^\prime = 0}}^{\substack{r}_{t-1}} \mathcal{N} (\boldsymbol{x}_t | \boldsymbol{\mu}_{r^\prime}, \Lambda_{r^\prime}) \cdot (1 - \frac{1}{\lambda}) \nonumber \\
            &= \frac{\lambda - 1}{\lambda} \cdot \prod_{\substack{r^\prime = 0}}^{\substack{r}_{t-1}} \mathcal{N} (\boldsymbol{x}_t | \boldsymbol{\mu}_{r^\prime}, \Lambda_{r^\prime}) \cdot P(r_{t-1}, \boldsymbol{x}_{0:t-1}). \nonumber
        \end{align*}

        Therefore, the permanence probability at time $t$, given that we consider the time of the previous changepoint $t_{cp}$ to be zero, is:
        \begin{align}
            P(&r_t = r_{t-1} + 1, \boldsymbol{x}_{0:t}) \nonumber\\ &= \frac{\lambda - 1}{\lambda} \cdot \prod_{\substack{r^\prime = 0}}^{\substack{r}_{t-1}} \mathcal{N} (\boldsymbol{x}_t | \boldsymbol{\mu}_{r^\prime}, \Lambda_{r^\prime}) \cdot P(r_{t-1}, \boldsymbol{x}_{0:t-1}) \nonumber \\
            &= \frac{\lambda - 1}{\lambda} P(r_{t-1}, \boldsymbol{x}_{t_{cp}:t-1})  
                \mathcal{N} (\boldsymbol{x}_t | \boldsymbol{\mu}_{t-1}^{(r)}, \Lambda_{t-1}^{(r)}). \nonumber
        \end{align}

        \item \textit{\textbf{Changepoint Probability:}} It follows from the results above and the fact that the changepoint probability is defined in Eq. \eqref{eq:changepoint-prob-def} as $P(r_t=0, \boldsymbol{x}_{1:t}) = \sum_{\substack{r_{t-1}}} P(r_{t-1}, \boldsymbol{x}_{1:t-1}) \pi_t^{(r)} H(r_{t-1})$, that the probability of jumping to a new situation is:
        \begin{align*}
            P(&r_t = 0, \boldsymbol{x}_{t_{cp}:t}) \\ &= \frac{1}{\lambda}  \cdot \prod_{\substack{r^\prime = 0}}^{\substack{r}_{t-1}} \mathcal{N} (\boldsymbol{x}_t | \boldsymbol{\mu}_{r^\prime}, \Lambda_{r^\prime}) \cdot   \sum_{\substack{\tau = 0}}^{r_{t-1}} P(\tau, \boldsymbol{x}_{t_{cp}:t-1})  \\
            &= \frac{1}{\lambda} \sum_{\substack{\tau = 0}}^{r_{t-1}} P(\tau, \boldsymbol{x}_{t_{cp}:t-1}) 
                \mathcal{N} (\boldsymbol{x}_t | \boldsymbol{\mu}_{\tau}^{(r)}, \Lambda_{\tau}^{(r)}).
        \end{align*}
    \qed
    \end{itemize}
\end{proof}

\end{document}
